\section{相平面：把解画成几何}
\label{sec:15-Differential-Equations-and-Dynamical-Systems/02-Stability-and-Dynamical-Systems/01}

一个两变量的库存与价格模型摆在你面前，右端写得出来，解写不出来。会上真正被问的却只有一句：这套机制放着不管，最后会稳在什么地方，还是会一直来回摆？

这个问题不需要解方程。把状态空间画出来，标出几个不动点，看清附近的箭头往哪转，答案就在图上。本节把上一单元的求解视角换成几何视角：放弃``\(t\) 时刻在哪里''，只保留``经过哪里、去向哪里''。

\subsection{自治系统在平面上钉出一张向量场}

二维自治系统写成

\[
\dot x=f(x,y),
\qquad
\dot y=g(x,y),
\]

或紧凑地 \(\dot u=F(u)\)，其中 \(u=(x,y)^{\trans}\)，\(F=(f,g)^{\trans}\)。自治指 \(F\) 不显含 \(t\)：规则不随时间改变，平面上每一点的运动方向是钉死的。整个平面因此是一张贴好的向量场，每个状态点上都插着一支箭头，说明现在若在这里，下一瞬往哪走。

轨迹是处处与向量场相切的曲线。上一单元\hyperref[sec:15-Differential-Equations-and-Dynamical-Systems/01-ODE-Theory/03]{``存在唯一性为什么需要 Lipschitz''}中的定理在这里兑现成一句几何断言：过每个点有且只有一条轨迹，因而轨迹互不相交，平面被梳理成一族流线。没有这块地基，相平面上就是一团乱麻，本节后面所有``读图''的操作都不合法。存在唯一性的价值，换成几何眼光才看得清楚。

\subsection{不动点是长期行为的第一批候选终点}

若 \(F(u^\ast)=0\)，则常数函数 \(u(t)\equiv u^\ast\) 是解，系统在那里永远停住。这样的点叫不动点，也叫平衡点。轨迹能做的事情不多：停在某个不动点上、绕闭轨周期运动、或者走向无穷，所以先找不动点总是对的。

在不动点附近令 \(u=u^\ast+w\)，把 \(F\) 展开，

\[
\dot w=F(u^\ast+w)=F(u^\ast)+Jw+o(\norm{w})=Jw+o(\norm{w}),
\]

其中 \(J\) 是 \(F\) 在 \(u^\ast\) 处的 Jacobi 矩阵。扔掉高阶项就得到线性化系统 \(\dot w=Jw\)，而\hyperref[sec:15-Differential-Equations-and-Dynamical-Systems/01-ODE-Theory/02]{``线性系统与矩阵指数''}已经把这类系统的行为全部交给了特征值。于是问``非线性系统在不动点附近怎样运动''，第一步永远是查 \(J\) 的谱。

\subsection{四种典型相图撑起全部分类}

\begin{center}
\begin{tikzpicture}[>=stealth]
  \begin{scope}
    \foreach \a in {0,45,90,135,180,225,270,315}{
      \draw[->] ({1.05*cos(\a)},{1.05*sin(\a)}) -- ({0.18*cos(\a)},{0.18*sin(\a)});
    }
    \fill (0,0) circle (1.6pt);
    \node at (0,-1.45) {\small 稳定结点};
  \end{scope}

  \begin{scope}[shift={(3.2,0)}]
    \draw[->] (-1.1,0) -- (-0.16,0);
    \draw[->] (1.1,0) -- (0.16,0);
    \draw[->] (0,0.16) -- (0,1.1);
    \draw[->] (0,-0.16) -- (0,-1.1);
    \foreach \sx/\sy in {1/1,1/-1,-1/1,-1/-1}{
      \draw[->] plot[smooth] coordinates
        {({\sx*1.05},{\sy*0.12}) ({\sx*0.6},{\sy*0.2}) ({\sx*0.4},{\sy*0.3})
         ({\sx*0.3},{\sy*0.4}) ({\sx*0.2},{\sy*0.6}) ({\sx*0.16},{\sy*0.78})};
    }
    \fill (0,0) circle (1.6pt);
    \node at (0,-1.45) {\small 鞍点};
  \end{scope}

  \begin{scope}[shift={(0,-3.3)}]
    \draw plot[domain=0:15.7,samples=90,variable=\t]
      ({1.05*exp(-0.09*\t)*cos(\t r)},{1.05*exp(-0.09*\t)*sin(\t r)});
    \draw[->] (1.05,0) -- (1.03,0.12);
    \fill (0,0) circle (1.6pt);
    \node at (0,-1.45) {\small 稳定螺旋};
  \end{scope}

  \begin{scope}[shift={(3.2,-3.3)}]
    \draw (0,0) circle (0.4);
    \draw (0,0) circle (0.72);
    \draw (0,0) circle (1.04);
    \draw[->] (0.4,-0.06) -- (0.4,0.06);
    \draw[->] (0.72,-0.06) -- (0.72,0.06);
    \draw[->] (1.04,-0.06) -- (1.04,0.06);
    \fill (0,0) circle (1.6pt);
    \node at (0,-1.45) {\small 中心};
  \end{scope}
\end{tikzpicture}
\end{center}

\emph{Jacobi 矩阵的特征值决定不动点附近长什么样；把箭头方向全部反过来，就得到对应的不稳定情形。}

图上四幅小相图对应四类谱。两个负实特征值时各方向一同衰减，轨迹几乎沿直线沉入不动点，是稳定结点；两个都为正则箭头反向，成为不稳定结点。实特征值一正一负给出鞍点：沿一个方向被吸入，沿另一个方向被抛出，绝大多数轨迹先靠近再被弹开。复共轭对 \(\alpha\pm i\beta\) 带来旋转，\(\alpha<0\) 是稳定螺旋，\(\alpha>0\) 是不稳定螺旋。纯虚特征值给出中心：一整族绕不动点的闭轨，既不靠近也不远离。

前面几类有一个共同点，特征值实部都不为零，这样的不动点称为双曲的。双曲情形下上面的读法给出确定结论，而中心是临界情形，稍后单独处理。

\subsection{鞍点的两条特殊轨迹是相图的分水岭}

鞍点值得多看一眼。线性鞍点上有两条穿过不动点的直线轨迹：沿负特征值方向的进入线和沿正特征值方向的离开线。非线性系统里它们弯成两条与特征方向相切的曲线，分别叫稳定流形（沿它前进的轨迹最终汇入鞍点）与不稳定流形（把时间反过来才回到鞍点）。

稳定流形是整张相图的分水岭。它把平面切成若干区域，初值落在哪一侧，就决定了轨迹最终被弹向哪个归宿。回到开头那个库存与价格的模型：若图上出现一个鞍点，那条稳定流形就是``救得回来''与``救不回来''的分界线，而不是某个凭经验拍出来的阈值。画相图时先定出不动点、再画出鞍点的两条流形，骨架就立起来了。

\subsection{线性化在双曲点上可靠，在临界点上沉默}

扔掉高阶项到底合不合法？Hartman--Grobman 定理给的答案是：在双曲不动点附近，非线性流与它的线性化流拓扑等价，存在一个连续的坐标变形把一族轨迹映成另一族。只要特征值实部全不为零，上面那四幅图的分类对原系统逐字成立。双曲是常态，非双曲才是意外。

中心正是那个意外。线性化给出纯虚特征值时，任意小的非线性项都能推翻结论。看一对只差一个符号的系统：

\[
\dot x=-y+x(x^2+y^2),
\qquad
\dot y=x+y(x^2+y^2),
\]

与

\[
\dot x=-y-x(x^2+y^2),
\qquad
\dot y=x-y(x^2+y^2).
\]

两者在原点的 Jacobi 矩阵完全相同，特征值都是 \(\pm i\)，线性化都判为中心。但在极坐标下，前者满足 \(r\dot r=x\dot x+y\dot y=(x^2+y^2)^2\)，即 \(\dot r=r^3\)，轨迹螺旋飞出，原点实为不稳定螺旋；后者给出 \(\dot r=-r^3\)，轨迹螺旋沉入，原点实为稳定螺旋。同一张线性化图纸，配上两个正负相反的三次项，结论完全相反。要裁决这种情形，得换一件不依赖线性化的工具，那是下一节的事。

\subsection{极限环让振荡自己定下振幅}

中心的闭轨有一整个连续族，初值决定停在哪一圈，换个初值就换一圈。工程与自然界里的自持振荡不是这样：钟表、心跳、电子振荡器的周期和振幅由系统自身定下，受了扰动还会回到原来的节律。对应的几何对象是极限环——一条孤立的闭轨，邻近轨迹从两侧或一侧贴上去。

标准例子是 van der Pol 振子

\[
\ddot x-\mu(1-x^2)\dot x+x=0,
\qquad \mu>0.
\]

与谐振子相比，它的阻尼系数 \(-\mu(1-x^2)\) 随振幅变号：\(\abs{x}>1\) 时为正，系统耗能，振幅被压回；\(\abs{x}<1\) 时为负，系统注入能量，振幅被吹大。两种机制在某个振幅上打平，形成一条稳定的极限环，除原点外的任何初值最终都贴上去，按系统自定的振幅与周期运动，无须外力持续驱动。

相平面上于是出现了两类长期归宿：不动点与极限环。二维平面上``除此以外还能怎样''由 Poincar\'e--Bendixson 定理严格限定，本 Part 不展开；而维数一旦升到三，这条限制就失效，后面讲混沌时会看到失效的后果。眼下留下的缺口是那个临界情形——线性化说不出话的时候，稳定性该怎么判。下一节换一个视角：不看谱，看能量。
